\section{Diagnostic causal reasoning: formal calculations}
\label{annex:p19-rat73-calculations}

\begin{annexitem}{annex:p19.rat73-elementary}{thm:p19-rat73-elementary}{Bayes, noisy-OR power and production attribution}
\begin{proof}
By the definition of conditional probability,
\(P(c\cap e)=P(e\mid c)P(c)=ub_c\).
The disjoint partition by \(c,\neg c\) gives
\(P(e)=ub_c+v(1-b_c)>0\). Division gives the first formula.

In the additional construction, \(E=1\) precisely when at least one
of \(C\land X\) or \(Z\) is one. Conditional on \(C=1\),
independence makes the probability of their joint failure
\((1-w_c)(1-w_a)\). Thus
\(u=1-(1-w_c)(1-w_a)=w_c+w_a-w_cw_a\).
Conditional on \(C=0\), only background production remains and
\(v=w_a\). Substituting these into total probability gives
\[
 P(e)=b_c(w_c+w_a-w_cw_a)+(1-b_c)w_a
     =b_cw_c+w_a-b_cw_cw_a.
\]
Substitution into the first formula gives the noisy-OR diagnostic
expression. Since \(u-v=w_c(1-w_a)\), division by \(1-v\)
when \(v<1\) identifies \(w_c\). Its nonnegative range forces
\(u\geq v\); \(u\leq1\) gives \(w_c\leq1\).
When \(v=1\), \(u=1\) for every \(w_c\), so the inversion
is genuinely unidentifiable from those conditional probabilities.

Let \(F=C\land X\), the event denoted \(c\to e\).
Independence gives \(P(F)=b_cw_c\), and \(F\subseteq e\).
Therefore \(P(F\mid e)=P(F)/P(e)\), which proves attribution.
Subtracting its numerator from the diagnostic numerator gives
\(b_cw_a(1-w_c)\geq0\), establishing the bound and difference.
If \(w_c=0\), \(P(e)=w_a>0\), attribution is zero, and
the diagnostic expression reduces to \(b_c\).
If \(w_a=0\), positive \(P(e)=b_cw_c\) makes both ratios one.

For the source's exact Sample~1 plug-in values,
\(u=3/10,v=1/10,b_c=1/2\), inversion gives \(w_c=2/9\).
Then \(P(e)=1/5\), \(ub_c/P(e)=(3/20)/(1/5)=3/4\),
and \(b_cw_c/P(e)=(1/9)/(1/5)=5/9\).
The numeric comparison thus follows without using a rounded power
as an exact input.
\end{proof}
\end{annexitem}

\begin{annexitem}{annex:p19.rat73-averaging}{thm:p19-rat73-averaging}{The conditioning order in an average of models}
\begin{proof}
The mixture law assigns an event \(H\) probability
\(\sum_i\rho_iP_i(H)\). Since
\(P_i(c\cap e)=q_it_i\), its joint numerator is
\(\sum_i\rho_iq_it_i\), and its effect denominator is
\(\sum_i\rho_iq_i>0\). Their ratio is the stated conditional.
If all \(q_i=q>0\), the common \(q\) cancels and leaves
\(\sum_i\rho_it_i\). Otherwise the effective conditional weight
of model \(i\) is \(\rho_iq_i/\sum_j\rho_jq_j\), not
generally \(\rho_i\).

For the first specified noisy-OR model,
\(P_0(e)=1/10\) and \(P_0(c\cap e)=1/20\), hence
\(t_0=1/2\). For the second,
\(u_1=4/5+1/5-4/25=21/25\),
\(q_1=(21/25)(1/2)+(1/5)(1/2)=13/25\), and
\(P_1(c\cap e)=21/50\), hence \(t_1=21/26\).
With equal weights, the average of the two conditional scores is
\((1/2+21/26)/2=17/26\).
The mixture's joint probability is
\((1/20+21/50)/2=47/200\), and its effect probability is
\((1/10+13/25)/2=31/100\). The conditional is consequently
\((47/200)/(31/100)=47/62\), not \(17/26\).

The same distinction occurs under a parameter distribution: an integral
of \(P_\theta(c\mid e)\) averages ratios, while conditioning the
integrated joint law divides the integrated joint numerator by the
integrated effect probability. The construction above uses a distribution
with two atoms, and therefore already supplies a counterexample to a
general equivalence. No claim about all continuously distributed models
is needed to establish that counterexample.
\end{proof}
\end{annexitem}

\begin{annexitem}{annex:p19.rat73-network}{thm:p19-rat73-network}{Factorizations, diagnostic odds and separating mechanisms}
\begin{proof}
Each displayed factorization is a nonnegative normalized mass function:
sum first over child variables using their normalized kernels, then
over their parents. For the common-effect mass, summing over \(b\)
gives \(p(a,c)=p(a)p(c)\), proving marginal independence.
For the common-cause mass and a value \(b\) of positive probability,
division by \(p(b)\) gives
\(p(a,c\mid b)=p(a\mid b)p(c\mid b)\).

For the chain, summing its joint mass over \(c\) gives
\(p(a,b)=p(a)p(b\mid a)\). Summing also over \(a\)
gives the marginal \(p(b)\). At a positive-probability \(b\),
\[
 p(a,c\mid b)=\frac{p(a)p(b\mid a)}{p(b)}p(c\mid b)
             =p(a\mid b)p(c\mid b).
\]
These divisions make no claim at a zero-probability conditioning value.

For a common-cause law with binary \(B\), positive prior odds and
positive likelihood denominators, Bayes' rule gives
\[
 \frac{P(b\mid a,c)}{P(\neg b\mid a,c)}
 =\frac{P(b)P(a,c\mid b)}{P(\neg b)P(a,c\mid\neg b)}.
\]
Applying the preceding conditional factorization in both worlds proves
the product-of-likelihood-ratios formula stated in the body.
It is essential to factor both likelihoods, not merely one of them.

To see why marginal dependence is not forced, choose the common-cause
kernels for \(A,C\) to ignore \(B\) and to be independent fair
binary outputs. Then \(p(a,c)=p(a)p(c)\) despite the same listed
graph. In the chain, choose \(p(c\mid b)=p(c)\), again giving
marginal independence of \(A,C\). A graph edge can be ineffective
in its kernel, so a faithfulness assumption is additional.

For the OR mechanism with independent binary causes,
\(P(b)=r+s-rs\), and \(a\subseteq b\) gives
\(P(a\mid b)=r/(r+s-rs)\).
Since \(0<r,s<1\), the denominator is positive and less than one,
so this value is strictly larger than \(r\).
When \(c\) is known true, \(b\) is automatically true; independence
then yields \(P(a\mid b,c)=P(a\mid c)=r\).
The logical-contrast example has \(r=1/10,s=2/5\), hence
\(P(b)=23/50\), \(P(a\mid b)=5/23\), and
\(P(c\mid b)=20/23\).

For the synergistic counterexample, every parent pair has probability
\(1/4\). The four positive-effect contributions are
\(1/40,1/20,1/20,1/4\), so \(P(b)=3/8\).
The two contributions with \(a\) true sum to \(3/10\),
and \(P(a\mid b)=(3/10)/(3/8)=4/5\).
The contributions with \(c\) true also sum to \(3/10\),
and \(P(a,b,c)=1/4\), so
\(P(a\mid b,c)=(1/4)/(3/10)=5/6\).
Every increase of either parent from absence to presence raises the
effect likelihood for both values of the other parent, but
\(5/6>4/5\). Thus a Markov collider, even with individually
productive influences, does not entail strict explaining away.
\end{proof}
\end{annexitem}
